\documentclass[10pt,a4paper]{amsart}
\usepackage[margin=19mm]{geometry}
\usepackage{amsmath,amssymb,mathtools}
\usepackage[scr=rsfs,cal=euler]{mathalfa}
\usepackage{xfrac}
\usepackage[unicode,hidelinks,bookmarks=false]{hyperref}
\newcommand{\ie}{i.e.\ }
\newcommand{\cf}{cf.\ }
\newcommand{\resp}{respectively\ }
\newcommand{\Subsection}{\subsection}
\providecommand{\nobreakdash}{\nobreak\hbox{-}}
\setlength{\emergencystretch}{2em}
\multlinegap0pt
\usepackage{amssymb}
\newtheorem{thm}{Theorem}
\def\geq{\geqslant}
\def\leq{\leqslant}
\title{{\bf \Large Local well-posedness for the quasilinear Schr\"odinger equations via the generalized energy method }}
\author{Jie Shao, Yi Zhou}
\date{Source: arXiv:2311.02556v3 (CC BY 4.0)}
\begin{document}
\maketitle

\noindent\emph{Source and licence.} {\bf \Large Local well-posedness for the quasilinear Schr\"odinger equations via the generalized energy method }. Version arXiv:2311.02556v3 (CC BY 4.0), distributed by arXiv under the Creative Commons Attribution 4.0 International licence, http://creativecommons.org/licenses/by/4.0/, as recorded in the arXiv licence metadata for this exact version and shown on the abstract page https://arxiv.org/abs/2311.02556v3. Full licence notice: \texttt{licenses/arxiv-2311.02556-CC-BY-4.0.txt}.\par\smallskip

\noindent\textit{Editorial scope.} Counted unit: the article's thm thm3 (source0.tex lines 235--238). The article's complete original proof of it is delivered with it (source0.tex lines 2408--2477, one \texttt{proof} environment of 70 lines, which the article places away from the statement). No mathematics is written, repaired or re-derived: the statement and the proof are the article's own lines, in the article's own notation, and no retained line is substituted or re-flowed.  Retained uncounted as the source-local context the proof needs: lines 103--108; lines 126--131; lines 142--147; lines 227--234; lines 435--441; lines 921--926; lines 928--930; lines 1037--1041; lines 2203--2211; lines 2295--2298; lines 2312--2316; lines 2388--2390.  Nothing was omitted from any retained span, so the package carries no omission record.  No retained line contains a citation, so the wrapper carries no bibliography and no bibliography entry.  No retained line contains a figure environment, an image-inclusion command or drawing code, so no image is delivered and the package declares no asset dependency.  Editorial formatting only, in addition: an AMS \texttt{amsart} preamble that restates the article's own declarations and macros used by the retained lines, namely the environments \texttt{thm} on separate counters, together with the standard packages those lines need.  The retained environments print in this document as Theorem 1 in order of appearance, whereas the article prints the same items with the numbering of its own full document.  The complete native source of the article as distributed in the arXiv e-print for this exact version is bundled unchanged as \texttt{sources/149547/source0.tex}.
	\begin{equation}	 \label{qNLS}
	\begin{cases}
			\sqrt{-1}\phi_t+\partial_i\left( g^{ij}\left(\phi,\overline{\phi}\right) \partial_j \phi\right)=F\left(\phi,\overline{\phi},\nabla\phi, \nabla \overline{\phi}\right),~~~\phi: \mathbb{R}^d \times \mathbb{R}^d\rightarrow \mathbb{C}^m\\
			\phi(0,x)=\phi_0(x),
	\end{cases}
	\end{equation}

\begin{align} \label{g0}
	\left(g_0^{ij}\right)=  \left(\begin{array}{cc}
	I_l &0\\
		0 & -I_{d-l}	
	\end{array}\right),
\end{align}

	\begin{equation}	 \label{qNLS0}
	\begin{cases}
		\sqrt{-1}\phi_t+g^{ij}\left(\phi,\overline{\phi},\nabla\phi, \nabla \overline{\phi} \right) \partial_{ij} \phi=F\left(\phi,\overline{\phi},\nabla\phi, \nabla \overline{\phi}\right),~~~\phi: \mathbb{R}^d \times \mathbb{R}^d\rightarrow \mathbb{C}^m\\
		\phi(0,x)=\phi_0(x).
	\end{cases}
\end{equation}

\begin{thm} \label{thm2}
$ 	(a) $ Suppose that    $ s>\frac {d}2+3, d \geq 1  $,  \eqref{qNLS} is  quadratic interaction problem and
\begin{align*}
	   \|\phi_0\|_{H^s(\mathbb{R}^d)} +\left\| \langle x\rangle^N \phi_0 \right\|_{L^2(\mathbb{R}^d)}\ll 1
\end{align*}  
  for   some integer $ N $ that depends only on $ d, s $.   Then  the ultrahyperbolic equation \eqref{qNLS} under condition   \eqref{g0} is locally wellposed in $   H^{s} (\mathbb{R}^d) \cap L^2(\langle x\rangle^N \mathrm{d}x) $ for $ t\in [0,1] $.\\
$ (b) $ The same results holds for \eqref{qNLS0} with $ s>\frac{d}{2}+4, d\geq 1. $
\end{thm}

\begin{align} \label{aconsm}
	&	-\partial_t \int_{\mathbb{R}^{d-1}}  \mathrm{Im}\left( \phi_{\alpha } \partial_k \overline{\phi}_\alpha \right) \mathrm{d} \hat{x}_k-\partial_{kk} \int_{\mathbb{R}^{d-1}} \mathrm{Re}\left( \phi_{\alpha }   g^{kj} \overline{\phi}_{\alpha j}\right) \mathrm{d} \hat{x}_k\nonumber\\
	&+2 \partial_k \int_{\mathbb{R}^{d-1}} \mathrm{Re}\left(\partial_k \phi_{\alpha }  g^{kj} \overline{\phi}_{\alpha j}\right) \mathrm{d} \hat{x}_k + \int_{\mathbb{R}^{d-1}}  \mathrm{Re}\left(\partial_{i} \phi_{\alpha }  \partial_k h^{ij} \partial_j \overline{\phi}_{\alpha }\right)  \mathrm{d} \hat{x}_k \\
	=&\,- \partial_k\int_{\mathbb{R}^{d-1}}   \mathrm{Re}\left( \overline{\phi }_{\alpha } \mathcal{N}\right)\mathrm{d} \hat{x}_k+2\int_{\mathbb{R}^{d-1}} \mathrm{Re} \left(\mathcal{N} \partial_k \overline{\phi}_\alpha\right)   \mathrm{d} \hat{x}_k,\nonumber\\
		&-\sum_{i \neq k} \varepsilon \partial_k \partial_k\partial_k  \int_{\mathbb{R}^{d-1}} \mathrm{Im}\left(\partial_{kk} \phi_{\alpha }  \overline{\phi}_\alpha\right)  \mathrm{d} \hat{x}_k+2 \varepsilon \partial_{k}\partial_{k}\int_{\mathbb{R}^{d-1}} \mathrm{Im}\left(\partial_{kj} \phi_{\alpha } \partial_j \overline{\phi}_\alpha\right) \mathrm{d} \hat{x}_k \nonumber\\
	&\,	+2 \varepsilon \partial_{k}\partial_{k}\int_{\mathbb{R}^{d-1}} \mathrm{Im}\left(\partial_{kk} \phi_{\alpha } \partial_k \overline{\phi}_\alpha\right) \mathrm{d} \hat{x}_k-2\varepsilon\mathrm{Im}\int_{\mathbb{R}^{d-1}} \left(\partial_{i}\partial_{jj} \phi_\alpha\partial_i\partial_k \overline{\phi}_\alpha\right)  \mathrm{d} \hat{x}_k,\nonumber
\end{align} 

	\begin{equation}	 \label{aqNLS}
		\begin{cases}
			\sqrt{-1}\phi_t+	\varepsilon\sqrt{-1}\Delta^2 \phi+\partial_i\left( g^{ij}\left(\phi,\overline{\phi}\right) \partial_j \phi\right)=F\left(\phi,\overline{\phi},\nabla\phi, \nabla \overline{\phi}\right)	,~~~\phi: t \times \mathbb{R}^d\rightarrow \mathbb{C}^m\\
			\phi(0,x)=\phi_0(x).
		\end{cases}
	\end{equation}

	\begin{align} \label{aqNLS1}
		\sqrt{-1}\phi_{\alpha t}+	\varepsilon\sqrt{-1}\Delta^2 \phi_\alpha+\partial_i\left( g^{ij} \partial_j \phi_{\alpha }\right)=F_z\nabla \phi_\alpha+F_{\overline{z}}\nabla \overline{\phi}_\alpha+\partial h^{ij}\partial_{ij } \phi_{\alpha -1}+G.
	\end{align}

 	\begin{align} \label{consm1-2}
 		&	-\partial_t \int_{\mathbb{R}^{d-1}}  \mathrm{Im}\left(v \partial_k \overline{v} \right) \mathrm{d} \hat{x}_k-\partial_{kk} \int_{\mathbb{R}^{d-1}} \mathrm{Re}\left( v   g^{kj} \left(\phi_1,\overline{\phi}_1\right) \overline{v}_{ j}\right) \mathrm{d} \hat{x}\nonumber\\
 		&+2 \partial_k \int_{\mathbb{R}^{d-1}} \mathrm{Re}\left(\partial_k v  g^{kj} \left(\phi_1,\overline{\phi}_1\right) \overline{v}_{  j}\right) \mathrm{d} \hat{x} + \int_{\mathbb{R}^{d-1}}  \mathrm{Re}\left(\partial_{i} v \partial_k h^{ij} \left(\phi_1,\overline{\phi}_1\right) \partial_j \overline{v} \right)  \mathrm{d} \hat{x} \\
 		=&\,- \partial_k\int_{\mathbb{R}^{d-1}}   \mathrm{Re}\left( \overline{v}  N_{uni} \right)\mathrm{d} \hat{x}+2\int_{\mathbb{R}^{d-1}} \mathrm{Re} \left( N_{uni} \partial_k \overline{ v} \right)   \mathrm{d} \hat{x}. \nonumber
 	\end{align}

\begin{align} \label{cener_c1_4_1}
	&\int_{\mathbb{R}}  \int_{\mathbb{R}^{d-1}}  \mathrm{Im}\left( \phi_{\alpha } \partial_k \overline{\phi}_\alpha \right) \mathrm{d} \hat{x}_k \int_{-\infty}^{x_k} \|\Lambda_k^{\frac 12}\phi_{\beta}(x_k,\cdot)\|^2_{L^2(\mathbb{R}^{d-1})} \mathrm{d} y_k  \mathrm{d} x_k \nonumber\\
	\lesssim &\,\left\|   D_k^{\frac 12}\phi_\alpha\right\|_{L^2(\mathbb{R}^d)}^2 \int_{-\infty}^{\infty} \|\Lambda_k^{\frac 12}\phi_{\beta}(x_k,\cdot)\|^2_{L^2(\mathbb{R}^{d-1})}\mathrm{d} x_k\nonumber\\
	&+\left\|   D_k^{\frac 12}\phi_{\alpha}\right\|_{L^2(\mathbb{R}^d)} \left\|   \phi_{\alpha}\right\|_{L^2(\mathbb{R}^d)} \left\| \|\Lambda_k^{\frac 12}\phi_{\beta}(x_k,\cdot)\|^2_{L^2(\mathbb{R}^{d-1})}\right\|_{L^2(\mathbb{R}_{x_k})}\\
	\lesssim &\,\left\|   D_k^{\frac 12}\phi_\alpha\right\|_{L^2(\mathbb{R}^d)}^2 \|\phi\|_{H^{|\alpha|}(\mathbb{R}^d)}^2\nonumber\\
	&+\left\|   D_k^{\frac 12}\phi_{\alpha}\right\|_{L^2(\mathbb{R}^d)} \left\|   \phi_{\alpha}\right\|_{L^2(\mathbb{R}^d)}\left\| \|\Lambda_k^{\frac 12}\phi_{\beta}(\cdot,\hat{x}_k)\|_{L^4(\mathbb{R}_{x_k})} \right\|_{L^{2}(\mathbb{R}^{d-1})}^2\nonumber\\
	\lesssim &\,\left\|   D_k^{\frac 12}\phi_\alpha\right\|_{L^2(\mathbb{R}^d)}^2 \|\phi\|_{H^{|\alpha|}(\mathbb{R}^d)}^2\nonumber+\left\|   D_k^{\frac 12}\phi_{\alpha}\right\|_{L^2(\mathbb{R}^d)} \left\|   \phi_{\alpha}\right\|_{L^2(\mathbb{R}^d)}\left\| \|\Lambda_k^{\frac 12}\phi_{\beta}(\cdot,\hat{x}_k)\|_{H^{\frac14}(\mathbb{R}_{x_k})} \right\|_{L^{2}(\mathbb{R}^{d-1})}^2\nonumber\\
	\lesssim &\,\|\phi\|^4_{H^{s_3+\frac 12}(\mathbb{R}^{d})},\nonumber
\end{align}

\begin{align}  \label{cenerM}
	W(t) 
	\lesssim &\,  \|\phi_0\|^4_{H^{s_3+\frac 12}(\mathbb{R}^{d})} +  (1+t)\cdot\left(1+\|\phi\|_{H^{s_3+\frac 12}}^2\right)\|\phi\|_{H^{s_3+\frac 12}}^4+\left\|  \phi\right\|_{H^{s_3+\frac 12}(\mathbb{R}^d)}^2 \cdot W (t),
\end{align}

\begin{align} \label{cener1}
	\| \phi\|_{H^{ s_3 }(\mathbb{R}^d)}^2\lesssim &\, 		\| \phi_0\|_{H^{s_3 }(\mathbb{R}^d)}^2\nonumber\\
	&+\sum_{|\alpha|=0}^{s_3}\int_0^t\int_{\mathbb{R}^d} \left| \mathrm{Im} \left[  \left(F_z\nabla \phi_\alpha+F_{\overline{z}}\nabla \overline{\phi}_\alpha+\partial h^{ij}\partial_{ij } \phi_{\alpha -1}+G\right)\overline{\phi}_{\alpha }\right]\right| \mathrm{d} x\mathrm{d} s\\
	\lesssim &\,\| \phi_0\|_{H^{s_3 }(\mathbb{R}^d)}^2+W(t)+t\cdot\|\phi\|_{H^{|\alpha|+\frac12}(\mathbb{R}^{d})}^4. \nonumber
\end{align}

\begin{align}\label{cener3}
	&	\|\phi\|_{H^{s_3+\frac12}(\mathbb{R}^{d})}^2\lesssim 	\|\phi_0\|_{H^{s_3+\frac12}(\mathbb{R}^{d})}^2+ W(t)+t\cdot\|\phi\|_{H^{|\alpha|+\frac12}(\mathbb{R}^{d})}^4.
\end{align}

\begin{thm} \label{thm3}
	$ 	(a) $   Suppose that $  s>\frac {d+3}2, d \geq 2,   $   \eqref{qNLS} is  cubic interaction problem and $ \|\phi_0\|_{H^s(\mathbb{R}^d)}\ll 1 $. Then the ultrahyperbolic equation \eqref{qNLS} under condition   \eqref{g0}  is locally wellposed in $ H^s(\mathbb{R}^d) $  for $t\in  [0,1] $.\\
	$ (b) $ The same results holds for \eqref{qNLS0} with $ s>\frac{d+5}{2}, d\geq 2. $
\end{thm}

\begin{proof}[\textbf{Proof of Theorem \ref{thm3}}]
As in the proof of Theorem~\ref{thm2}, we employ the artificial viscosity method with the approximate equation \eqref{aqNLS}. Recall \eqref{aqNLS1}:
	\begin{align*}  
		\sqrt{-1}\phi_{\alpha t}+	\varepsilon\sqrt{-1}\Delta^2 \phi_\alpha+\partial_i\left( g^{ij} \partial_j \phi_{\alpha }\right)=F_z\nabla \phi_\alpha+F_{\overline{z}}\nabla \overline{\phi}_\alpha+\partial h^{ij}\partial_{ij } \phi_{\alpha -1}+G.
	\end{align*}
Analogous to to \eqref{cener1} and \eqref{cener3},
we derive the energy estimates
	\begin{align*}
		&	\| \phi\|_{H^{s_3 }(\mathbb{R}^d)}^2+\varepsilon  \sum_{i,j=1}^d\int_0^t\left\| \partial_{ij}\phi\right\|_{H^{s_3}(\mathbb{R}^d)}^2 \mathrm{d} s\\
		\lesssim &\, 		\| \phi_0\|_{H^{s_3 }(\mathbb{R}^d)}^2+\sum_{|\alpha|=0}^{s_3}\int_0^t\int_{\mathbb{R}^d} \left| \mathrm{Im} \left[  \left(F_z\nabla \phi_\alpha+F_{\overline{z}}\nabla \overline{\phi}_\alpha+\partial h^{ij}\partial_{ij } \phi_{\alpha -1}+G\right)\overline{\phi}_{\alpha }\right]\right| \mathrm{d} x\mathrm{d} s\\
		\lesssim &\,\| \phi_0\|_{H^{s_3 }(\mathbb{R}^d)}^2+W(t)+t\cdot\|\phi\|_{H^{s_3+\frac12}(\mathbb{R}^{d})}^6 \nonumber
	\end{align*}
	and 
	\begin{align} \label{acener3}
		\|\phi\|_{H^{s_3+\frac12}(\mathbb{R}^{d})}^2+\varepsilon  \sum_{i,j=1}^d\int_0^t\left\| \partial_{ij}\phi\right\|_{H^{s_3+\frac 12}(\mathbb{R}^d)}^2 \mathrm{d} s\lesssim 	\|\phi_0\|_{H^{s_3+\frac12}(\mathbb{R}^{d})}^2+ W(t)+t\cdot\|\phi\|_{H^{s_3+\frac12}(\mathbb{R}^{d})}^6.
	\end{align}
	
	
	
	Similar to \eqref{cener_c1_4_1}, we verify
	\begin{align*}
		&\varepsilon \left| \int_0^t	\int_{\mathbb{R}}\int_{-\infty}^{x_k} \|\Lambda_k^{\frac 12}\phi_{\beta}(y_k,\cdot)\|^2_{L^2(\mathbb{R}^{d-1})} \mathrm{d} y_k   \cdot\mathrm{Im}\int_{\mathbb{R}^{d-1}} \left(\partial_{kjj}\phi_{\alpha } \partial_{kk}\overline{\phi}_\alpha\right) \mathrm{d} \hat{x}_k\mathrm{d} x_k\mathrm{d} s\right|\\
		\lesssim &\,\varepsilon \int_0^t \left\|   D_k^{\frac 12}  \partial_{kk}\phi_\alpha\right\|_{L^2(\mathbb{R}^d)}  \left\|   D_k^{\frac 12} \partial_{jj}\phi_\alpha\right\|_{L^2(\mathbb{R}^d)} \|\Lambda_k^{\frac 12}\phi_{\beta}\|_{L^{2}(\mathbb{R}^d)}^2 \mathrm{d} s\nonumber\\
		&+\varepsilon\int_0^t \left\|    \partial_{kk}\phi_{\alpha}\right\|_{L^2(\mathbb{R}^d)}   \left\|   D_k^{\frac 12} \partial_{jj}\phi_\alpha\right\|_{L^2(\mathbb{R}^d)} \left\| \|\Lambda_k^{\frac 12}\phi_{\beta}(\cdot,\hat{x}_k)\|_{H^{\frac12}(\mathbb{R}_{x_k})} \right\|_{L^{2}(\mathbb{R}^{d-1})}^2 \mathrm{d} s\nonumber\\
		\lesssim &\,\varepsilon \left\| \Lambda_k^{\frac 12}\phi_{\beta}\right\|_{H^{\frac12}(\mathbb{R}^d)}^2  \sum_{i,j=1}^d\int_0^t\left\| \partial_{ij}\phi\right\|_{H^{s_3+\frac 12}(\mathbb{R}^d)}^2 \mathrm{d} s.
	\end{align*}
An analysis of \eqref{aconsm} analogous to \eqref{cenerM} then gives
	\begin{align*}
		&W(t) \lesssim   \|\phi_0\|^4_{H^{s_3+\frac 12}(\mathbb{R}^{d})} +  (1+t)\cdot\left(1+\|\phi\|_{H^{s_3+\frac 12}}^2\right)\|\phi\|_{H^{s_3+\frac 12}}^4+\left\|  \phi\right\|_{H^{s_3+\frac 12}(\mathbb{R}^d)}^2 \cdot W (t)\\
		&+\varepsilon \left\| \Lambda_k^{\frac 12}\phi_{\beta}\right\|_{H^{\frac12}(\mathbb{R}^d)}^2\int_0^t\left\| \nabla\phi\right\|_{H^{s_3+\frac 12}}^2 \mathrm{d} s. \nonumber
	\end{align*}
Combining this with \eqref{acener3}, we obtain
	\begin{align*}
		&W(t) +\|\phi\|_{H^{s_3+\frac12}(\mathbb{R}^{d})}^2+\varepsilon  \sum_{i,j=1}^d\int_0^t\left\| \partial_{ij}\phi\right\|_{H^{s_3+\frac 12}(\mathbb{R}^d)}^2 \mathrm{d} s\lesssim  \|\phi_0\|_{H^{s_3+\frac12}(\mathbb{R}^{d})}^2+\|\phi_0\|_{H^{s_3+\frac12}(\mathbb{R}^{d})}^4\nonumber\\
		&+(1+t)\cdot\left(1+\|\phi\|_{H^{s_3+\frac 12}}^2\right)\|\phi\|_{H^{s_3+\frac 12}}^4+\left\|  \phi\right\|_{H^{s_3+\frac 12}(\mathbb{R}^d)}^2 \cdot W (t)\nonumber\\
		&+\varepsilon  \left\|  \phi\right\|_{H^{s_3+\frac 12}(\mathbb{R}^d)}^2 \int_0^t\left\| \nabla\phi\right\|_{H^{s_3+\frac 12}}^2 \mathrm{d} s.
	\end{align*}
	For $ \|\phi_0\|_{H^{s_3+\frac12}(\mathbb{R}^{d})} $ small enough, we conclude that for \(t \in [0,1]\) and \(\frac{d+3}{2} < s_3 + \frac12\),
	\begin{align*} 
		W(t) +\|\phi\|_{H^{s_3+\frac12}(\mathbb{R}^{d})}^2+\varepsilon  \sum_{i,j=1}^d\int_0^t\left\| \partial_{ij}\phi\right\|_{H^{s_3+\frac 12}(\mathbb{R}^d)}^2 \mathrm{d} s \leq C\left(\|\phi_0\|_{H^{s_3+\frac12}(\mathbb{R}^{d})}\right).
	\end{align*}
	The existence of a solution follows by an argument similar to that in the proof of Theorem~\ref{thm2}.
	 
	 The proof of the uniqueness and continuity of the solution  is also similar with the quadratic interaction problem except the following points.  The weight to multiply by \eqref{consm1-2} is still
 $$ \int_{-\infty}^{x_k} \|\Lambda_k^{\frac 12}v_{\beta}(y_k,\cdot)\|^2_{L^2(\mathbb{R}^{d-1})} \mathrm{d} y_k  $$ and the space time norms for $ v(t) =\phi_1 -\phi_2$ is
	 \begin{align*}
	 	&W_v(t)	
	 	:=\sum_{k=1}^d\sum_{|\beta|\leq s_3-1} \int_0^t \int_{\mathbb{R}}  \| \Lambda_k^{\frac 12}v_{\beta}(x_k,\cdot)\|_{L^2(\mathbb{R}^{d-1})}^2  \|\partial_k v (x_k,\cdot)\|_{L^2(\mathbb{R}^{d-1})}^2 \mathrm{d} x_k \mathrm{d} s,
	 \end{align*}
	 and  
	 the solutions satisfy
	 \begin{align*}
	 	\label{csol1-2}
	 	&	\|\phi_l\|_{H^{s_3+\frac12}(\mathbb{R}^{d})}^2+	W(\phi_l)(t) \leq C\left(\|\phi_{l0}\|_{H^{s_3+\frac12}(\mathbb{R}^{d})}\right)\ll 1,~~l=1,2.
	 \end{align*}
	 Since 
	 \begin{align*}
	 	\max\{ \|h^{ij}_y\|_{L^{\infty}}, \|h^{ij}_{\overline{y}}\|_{L^{\infty}}\} \lesssim  \| \phi_1\|_{L^{\infty}}+\| \phi_2\|_{L^{\infty}},
	 \end{align*}
we have
	 \begin{align}
	 	&\int_0^t \int_{\mathbb{R}}\int_{-\infty}^{x_k} \|\Lambda_k^{\frac 12}v_{\beta}(y_k,\cdot)\|^2_{L^2(\mathbb{R}^{d-1})} \mathrm{d} y_k	\int_{\mathbb{R}^{d-1}} \mathrm{Re} \left(\partial_i\left( (g^{ij}\left(\phi_1,\overline{\phi}_1\right)-g^{ij}\left(\phi_2,\overline{\phi}_2\right)) \partial_j\phi_2\right)\partial_k \overline{v} \right)   \mathrm{d} \hat{x}  \mathrm{d} x_k \mathrm{d}s\nonumber\\
	 	=&\,\mathrm{Re}  \int_0^t \int_{\mathbb{R}} \int_{-\infty}^{x_k} \|\Lambda_k^{\frac 12}v_{\beta}(y_k,\cdot)\|^2_{L^2(\mathbb{R}^{d-1})} \mathrm{d} y_k	\int_{\mathbb{R}^{d-1}} \partial_i\left( h^{ij}_y  \left(\phi_1 -  \phi_2 \right)+h^{ij}_{\overline{y}} (\overline{\phi}_1-\overline{\phi}_2)\right) \partial_j\phi_2\partial_k \overline{v}  \mathrm{d} \hat{x}  \mathrm{d} x_k \mathrm{d}s \nonumber\\
	 	&+ \mathrm{Re}\int_0^t \int_{\mathbb{R}}  \int_{-\infty}^{x_k} \|\Lambda_k^{\frac 12}v_{\beta}(y_k,\cdot)\|^2_{L^2(\mathbb{R}^{d-1})} \mathrm{d} y_k	\int_{\mathbb{R}^{d-1}}   \left( h^{ij}_y  \left(\phi_1 -  \phi_2 \right)+h^{ij}_{\overline{y}} (\overline{\phi}_1-\overline{\phi}_2)\right)\partial_{ij}\phi_2\partial_k \overline{v}   \mathrm{d} \hat{x}  \mathrm{d} x_k \mathrm{d}s\nonumber \\
	 	\lesssim &\,\left( \left\|  \phi_1\right\|_{H^{s_3+\frac 12}(\mathbb{R}^d)}^2 + \left\|  \phi_2\right\|_{H^{s_3+\frac 12}(\mathbb{R}^d)}^2 \right) \cdot W_v(t)\\
	 	&+ \left( \left\|  \phi_1\right\|_{H^{s_3+\frac 12}(\mathbb{R}^d)}^2 + \left\|  \phi_2\right\|_{H^{s_3+\frac 12}(\mathbb{R}^d)}^2 \right) \cdot t^{\frac 12}\| \partial_{ij}\phi_2\|_{L^2(\mathbb{R}^d)} \|h^{ij}_y\|_{L^{\infty}(\mathbb{R}^d)} \nonumber\\
	 	&\cdot\left(\int_0^t \int_{\mathbb{R}}    \|\phi_1 -  \phi_2 (x_k,\cdot)\|_{L^\infty(\mathbb{R}^{d-1})}^2   \|\partial_k v (x_k,\cdot)\|_{L^2(\mathbb{R}^{d-1})}^2 \mathrm{d} x_k \mathrm{d} s  \right)\nonumber\\
	 	\lesssim &\,\left( \left\|  \phi_1\right\|_{H^{s_3+\frac 12}(\mathbb{R}^d)}^2 + \left\|  \phi_2\right\|_{H^{s_3+\frac 12}(\mathbb{R}^d)}^2 \right) \cdot W_v(t)+t^{\frac 12}\left( \left\|  \phi_1\right\|_{H^{s_3+\frac 12}(\mathbb{R}^d)}^2 + \left\|  \phi_2\right\|_{H^{s_3+\frac 12}(\mathbb{R}^d)}^2 \right)^2 \cdot W_v^{\frac 12}(t).\nonumber
	 \end{align}		
\end{proof}	

\end{document}
